\chapter{2019年天津卷（文）}

\section{单选题}

\begin{problem}
设集合 \(A=\{-1,1,2,3,5\}\)，\(B=\{2,3,4\}\)，\(C=\{x\in\R\mid 1\le x<3\}\)，则 \((A\cap C)\cup B=\)（\quad）

\choices
    {\(\{2\}\)}
    {\(\{2,3\}\)}
    {\(\{-1,2,3\}\)}
    {\(\{1,2,3,4\}\)}
\end{problem}

\begin{answer}
D
\end{answer}

\begin{solution}
因为 \(C=\{x\in\R\mid 1\le x<3\}\)，所以 \(A\cap C=\{1,2\}\). 于是 \((A\cap C)\cup B=\{1,2,3,4\}\). 故选 D.
\end{solution}

\begin{problem}
设变量 \(x\)，\(y\) 满足约束条件

\bitmapfigure[max width=4.8cm]{img/2019/tianjin_liberal/q2_fig1.png}

则目标函数 \(z=-4x+y\) 的最大值为（\quad）

\choices
    {\(2\)}
    {\(3\)}
    {\(5\)}
    {\(6\)}
\end{problem}

\begin{answer}
C
\end{answer}

\begin{solution}
由图可知可行域的顶点之一为 \(A(-1,1)\). 目标函数 \(z=-4x+y\) 在点 \(A\) 处取得最大值，\(z_{\max}=-4\times(-1)+1=5\). 故选 C.
\end{solution}

\begin{problem}
设 \(x\in\R\)，则“\(0<x<5\)”是“\(|x-1|<1\)”的（\quad）

\choices
    {充分而不必要条件}
    {必要而不充分条件}
    {充要条件}
    {既不充分也不必要条件}
\end{problem}

\begin{answer}
B
\end{answer}

\begin{solution}
\(|x-1|<1\) 等价于 \(0<x<2\). 故 \(0<x<5\) 推不出 \(|x-1|<1\)；由 \(|x-1|<1\) 能推出 \(0<x<5\). 故“\(0<x<5\)”是“\(|x-1|<1\)”的必要而不充分条件. 故选 B.
\end{solution}

\begin{problem}
阅读右边的程序框图，运行相应的程序，输出 \(S\) 的值为（\quad）

\texfigure[max width=6.0cm]{img_repaint/2019/tianjin_liberal/q4_fig1.tex}

\choices
    {\(5\)}
    {\(8\)}
    {\(24\)}
    {\(29\)}
\end{problem}

\begin{answer}
B
\end{answer}

\begin{solution}
按程序框图逐步计算：\(S=0\)，\(i=1\) 时 \(S=1\)；\(i=2\) 时 \(j=1\)，\(S=1+2\times 2=5\)；\(i=3\) 时 \(S=8\)，此时 \(i=4\)，满足 \(i\geq4\)，故输出 \(8\). 故选 B.
\end{solution}

\begin{problem}
已知 \(a=\log_2 7\)，\(b=\log_3 8\)，\(c=0.3^{0.2}\)，则 \(a\)，\(b\)，\(c\) 的大小关系为（\quad）

\choices
    {\(c<b<a\)}
    {\(a<b<c\)}
    {\(b<c<a\)}
    {\(c<a<b\)}
\end{problem}

\begin{answer}
A
\end{answer}

\begin{solution}
因为 \(0<0.3<1\)，所以 \(c=0.3^{0.2}<1\). 又 \(1<\log_3 8<\log_3 9=2\)，且 \(\log_2 7> \log_2 4=2\)，所以 \(c<b<a\). 故选 A.
\end{solution}

\begin{problem}
已知抛物线 \(y^2=4x\) 的焦点为 \(F\)，准线为 \(l\). 若与双曲线 \(\dfrac{x^2}{a^2}-\dfrac{y^2}{b^2}=1\)（\(a>0\)，\(b>0\)） 的两条渐近线分别交于点 \(A\) 和点 \(B\)，且 \(|AB|=4|OF|\)（\(O\) 为原点），则双曲线的离心率为（\quad）

\choices
    {\(\sqrt2\)}
    {\(\sqrt3\)}
    {\(2\)}
    {\(\sqrt5\)}
\end{problem}

\begin{answer}
D
\end{answer}

\begin{solution}
抛物线 \(y^2=4x\) 的焦点 \(F\) 为 \((1,0)\)，故 \(OF=1\)，准线 \(l\) 的方程为 \(x=-1\). 双曲线的渐近线方程为 \(y=\pm \dfrac{b}{a}x\)，故 \(l\) 与两条渐近线的交点为 \(A\!\left(-1,\dfrac{b}{a}\right)\)，\(B\!\left(-1,-\dfrac{b}{a}\right)\)，于是 \(|AB|=\dfrac{2b}{a}=4|OF|=4\)，从而 \(b=2a\). 故离心率 \(e=\sqrt{1+\left(\dfrac{b}{a}\right)^2}=\sqrt5\). 故选 D.
\end{solution}

\begin{problem}
已知函数 \(f(x)=A\sin(\omega x+\varphi)\)（\(A>0\)，\(\omega>0\)，\(|\varphi|<\pi\)） 是奇函数，且 \(f(x)\) 的最小正周期为 \(\pi\)，将 \(y=f(x)\) 的图象上所有点的横坐标伸长到原来的 2 倍（纵坐标不变），所得图象对应的函数为 \(g(x)\). 若 \(g\!\left(\dfrac{\pi}{4}\right)=\sqrt2\)，则 \(f\!\left(\dfrac{3\pi}{8}\right)=\)（\quad）

\choices
    {\(-2\)}
    {\(-\sqrt2\)}
    {\(\sqrt2\)}
    {\(2\)}
\end{problem}

\begin{answer}
C
\end{answer}

\begin{solution}
因为 \(f(x)\) 是奇函数，所以 \(\varphi=0\)，从而 \(f(x)=A\sin(\omega x)\). 将横坐标伸长到原来的 2 倍后，\(g(x)=A\sin\!\left(\dfrac{\omega x}{2}\right)\). 由 \(g(x)\) 的最小正周期为 \(2\pi\)，得 \(\omega=2\). 又 \(g\!\left(\dfrac{\pi}{4}\right)=A\sin\!\left(\dfrac{\pi}{4}\right)=\sqrt2\)，所以 \(A=2\). 因此 \(f\!\left(\dfrac{3\pi}{8}\right)=2\sin\!\left(\dfrac{3\pi}{4}\right)=\sqrt2\). 故选 C.
\end{solution}

\begin{problem}
已知函数 \(f(x)=2\sqrt{x}\)（\(0\le x\le 1\)），\(f(x)=\dfrac1x\)（\(x>1\)）. 若关于 \(x\) 的方程 \(f(x)=-\dfrac14x+a\)（\(a\in\R\)） 恰有两个互异的实数解，则 \(a\) 的取值范围为（\quad）

\choices
    {\(\left[\dfrac54,\dfrac94\right]\)}
    {\(\left(\dfrac54,\dfrac94\right]\)}
    {\(\left(\dfrac54,\dfrac94\right]\cup\{1\}\)}
    {\(\left[\dfrac54,\dfrac94\right]\cup\{1\}\)}
\end{problem}

\begin{answer}
D
\end{answer}

\begin{solution}
作出 \(f(x)\) 的图象与直线 \(y=-\dfrac14x+a\) 的图象，讨论交点个数可知：当直线与左支相交一次，与右支相交一次时，\(a\in\left[\dfrac54,\dfrac94\right]\)；当直线与右支相切时，\(a=1\). 故所求范围为 \(\left[\dfrac54,\dfrac94\right]\cup\{1\}\). 故选 D.
\end{solution}
\section{填空题}

\begin{problem}
\(i\) 是虚数单位，则 \(\left|\dfrac{5-i}{1+i}\right|=\fillinblank{}\).
\end{problem}

\begin{answer}
\(\sqrt{13}\)
\end{answer}

\begin{solution}
\(\left|\dfrac{5-i}{1+i}\right|=\dfrac{|5-i|}{|1+i|}=\sqrt{\dfrac{26}{2}}=\sqrt{13}\).
\end{solution}

\begin{problem}
设 \(x\in\R\)，使不等式 \(3x^2+x-2<0\) 成立的 \(x\) 的取值范围为\fillinblank{}.
\end{problem}

\begin{answer}
\(\left(-1,\dfrac23\right)\)
\end{answer}

\begin{solution}
\(3x^2+x-2<0 \iff (x+1)(3x-2)<0\)，故解集为 \(\left(-1,\dfrac23\right)\).
\end{solution}

\begin{problem}
曲线 \(y=\cos x-\dfrac x2\) 在点 \((0,1)\) 处的切线方程为\fillinblank{}.
\end{problem}

\begin{answer}
\(x+2y-2=0\)
\end{answer}

\begin{solution}
\(y'=-\sin x-\dfrac12\)，故在 \(x=0\) 处切线斜率为 \(-\dfrac12\). 切线方程为 \(y-1=-\dfrac12x\)，即 \(x+2y-2=0\).
\end{solution}

\begin{problem}
已知四棱锥的底面是边长为 \(\sqrt2\) 的正方形，侧棱长均为 \(\sqrt5\). 若圆柱的一个底面的圆周经过四棱锥四条侧棱的中点，另一个底面的圆心为四棱锥底面的中心，则该圆柱的体积为\fillinblank{}.
\end{problem}

\begin{answer}
\(\dfrac{\pi}{4}\)
\end{answer}

\begin{solution}
四棱锥的高为 \(\sqrt{5-1}=2\)，故圆柱的高为 \(1\)，圆柱的底面半径为 \(\dfrac12\). 因此 \(V=\pi\left(\dfrac12\right)^2\cdot 1=\dfrac{\pi}{4}\).
\end{solution}

\begin{problem}
设 \(x>0\)，\(y>0\)，\(x+2y=5\)，则 \(\dfrac{(x+1)(2y+1)}{xy}\) 的最小值为\fillinblank{}.
\end{problem}

\begin{answer}
\(\dfrac{98}{25}\)
\end{answer}

\begin{solution}
\(\dfrac{(x+1)(2y+1)}{xy}=\dfrac{2xy+x+2y+1}{xy}=2+\dfrac6{xy}\). 由 \(x+2y=5\) 可得 \(xy=x(5-x)/2\le \dfrac{25}{8}\)，当 \(x=\dfrac52\)，\(y=\dfrac54\) 时取等号. 故所求最小值为 \(2+\dfrac6{25/8}=\dfrac{98}{25}\).
\end{solution}

\begin{problem}
在四边形 \(ABCD\) 中，\(AD//BC\)，\(|AB|=2\sqrt3\)，\(|AD|=5\)，\(\angle A=30^\circ\)，点 \(E\) 在线段 \(CB\) 的延长线上，且 \(AE=BE\)，则 \(\overrightarrow{BD}\cdot\overrightarrow{AE}=\fillinblank{}\).
\end{problem}

\begin{answer}
\(-1\)
\end{answer}

\begin{solution}
建立直角坐标系，设 \(A(0,0)\)，\(D(5,0)\). 因为 \(|AB|=2\sqrt3\)，\(\angle A=30^\circ\)，所以 \(B(3,\sqrt3)\). 又因 \(AD//BC\) 且 \(E\) 在线段 \(CB\) 的延长线上，可设 \(E(x,\sqrt3)\). 由 \(AE=BE\) 知 \(E\) 在 \(AB\) 的垂直平分线上，故 \(E(1,\sqrt3)\). 于是 \(\overrightarrow{BD}=(2,-\sqrt3)\)，\(\overrightarrow{AE}=(1,\sqrt3)\)，从而 \(\overrightarrow{BD}\cdot\overrightarrow{AE}=2-3=-1\).
\end{solution}
\section{解答题}

\begin{problem}
2019 年，我国施行个人所得税专项附加扣除办法，涉及子女教育，继续教育，大病医疗，住房贷款利息或者住房租金，赡养老人等六项专项附加扣除. 某单位老，中，青员工分别有 \(72,108,120\) 人，现采用分层抽样的方法，从该单位上述员工中抽取 \(25\) 人调查专项附加扣除的享受情况.

\begin{enumerate}
\item 应从老，中，青员工中分别抽取多少人？
\item 抽取的 \(25\) 人中，享受至少两项专项附加扣除的员工有 \(6\) 人，分别记为 \(A\)，\(B\)，\(C\)，\(D\)，\(E\)，\(F\). 享受情况如右表，其中“\(\circ\)”表示享受，“\(\times\)”表示不享受. 现从这 \(6\) 人中随机抽取 \(2\) 人接受采访.
\begin{center}
\begin{tabular}{c|cccccc}
员工项目 & \(A\) & \(B\) & \(C\) & \(D\) & \(E\) & \(F\) \\
\hline
子女教育 & \(\circ\) & \(\circ\) & \(\times\) & \(\circ\) & \(\times\) & \(\circ\) \\
继续教育 & \(\times\) & \(\times\) & \(\circ\) & \(\times\) & \(\circ\) & \(\circ\) \\
大病医疗 & \(\times\) & \(\times\) & \(\times\) & \(\circ\) & \(\times\) & \(\times\) \\
住房贷款利息 & \(\circ\) & \(\circ\) & \(\times\) & \(\times\) & \(\circ\) & \(\circ\) \\
住房租金 & \(\times\) & \(\times\) & \(\circ\) & \(\times\) & \(\times\) & \(\times\) \\
赡养老人 & \(\circ\) & \(\circ\) & \(\times\) & \(\times\) & \(\times\) & \(\circ\)
\end{tabular}
\end{center}

\begin{enumerate}
\item 试用列举法写出所有可能的抽取结果；
\item 设事件 \(M\) 为“从这 \(6\) 人中随机抽取 \(2\) 人享受的专项附加扣除至少有一项相同”，求事件 \(M\) 发生的概率.
\end{enumerate}
\end{enumerate}
\end{problem}

\begin{solution}
\begin{enumerate}
\item 由已知，老，中，青员工人数之比为 \(72:108:120=6:9:10\). 又样本容量为 \(25\)，所以应从老，中，青员工中分别抽取 \(6\) 人，\(9\) 人，\(10\) 人.
\item
\begin{enumerate}
\item 从 \(A\)，\(B\)，\(C\)，\(D\)，\(E\)，\(F\) 中随机抽取 \(2\) 人的所有结果为
\(\{A,B\}\)，\(\{A,C\}\)，\(\{A,D\}\)，\(\{A,E\}\)，\(\{A,F\}\)，\(\{B,C\}\)，\(\{B,D\}\)，\(\{B,E\}\)，\(\{B,F\}\)，\(\{C,D\}\)，\(\{C,E\}\)，\(\{C,F\}\)，\(\{D,E\}\)，\(\{D,F\}\)，\(\{E,F\}\)，
共 \(15\) 种.

\item 设事件 \(M\) 为“从这 \(6\) 人中随机抽取 \(2\) 人享受的专项附加扣除至少有一项相同”，由表可知，满足“至少有一项相同”的结果有
\(\{A,B\}\)，\(\{A,D\}\)，\(\{A,E\}\)，\(\{A,F\}\)，\(\{B,D\}\)，\(\{B,E\}\)，\(\{B,F\}\)，\(\{C,E\}\)，\(\{C,F\}\)，\(\{D,F\}\)，\(\{E,F\}\)，
共 \(11\) 种，所以 \(P(M)=\frac{11}{15}\).
\end{enumerate}
\end{enumerate}
\end{solution}

\begin{problem}
在 \(\triangle ABC\) 中，内角 \(A\)，\(B\)，\(C\) 所对的边分别为 \(a\)，\(b\)，\(c\). 已知 \(b+c=2a\)，\(3c\sin B=4a\sin C\).

\begin{enumerate}
\item 求 \(\cos B\) 的值；
\item 求 \(\sin\left(2B+\dfrac{\pi}{6}\right)\) 的值.
\end{enumerate}
\end{problem}

\begin{solution}
\begin{enumerate}
\item 由正弦定理，\(\dfrac{b}{\sin B}=\dfrac{c}{\sin C}\)，故 \(b\sin C=c\sin B\). 又由 \(3c\sin B=4a\sin C\)，得 \(3b=4a\). 再由 \(b+c=2a\)，可得 \(b=\dfrac{4a}{3}\)，\(c=\dfrac{2a}{3}\). 由余弦定理可得 \(\cos B=\dfrac{a^2+c^2-b^2}{2ac}=-\dfrac14\).

\item 由（1）知 \(\sin B=\sqrt{1-\cos^2 B}=\dfrac{\sqrt{15}}{4}\)，从而 \(\sin 2B=2\sin B\cos B=-\dfrac{\sqrt{15}}{8}\)，\(\cos 2B=\cos^2 B-\sin^2 B=-\dfrac78\). 于是 \(\sin\left(2B+\dfrac{\pi}{6}\right)=\sin2B\cos\dfrac{\pi}{6}+\cos2B\sin\dfrac{\pi}{6}=-\dfrac{3\sqrt5+7}{16}\).
\end{enumerate}
\end{solution}

\begin{problem}
如图，在四棱锥 \(P-ABCD\) 中，底面 \(ABCD\) 为平行四边形，\(\triangle PCD\) 为等边三角形，平面 \(PAC\perp\) 平面 \(PCD\)，\(PA\perp CD\)，\(CD=2\)，\(AD=3\).

\begin{enumerate}
\item 设 \(G\)，\(H\) 分别为 \(PB\)，\(AC\) 的中点，求证：\(GH//\) 平面 \(PAD\)；
\item 求证：\(PA\perp\) 平面 \(PCD\)；
\item 求直线 \(AD\) 与平面 \(PAC\) 所成角的正弦值.
\end{enumerate}

\bitmapfigure[max width=8.5cm]{img/2019/tianjin_liberal/q17_fig1.png}
\end{problem}

\begin{solution}
\begin{enumerate}
\item 连接 \(BD\). 因为 \(ABCD\) 是平行四边形，所以 \(AC\) 与 \(BD\) 互相平分，设交点为 \(H\). 又 \(G\) 为 \(PB\) 的中点，所以在三角形 \(PBD\) 中，\(GH\) 是中位线，故 \(GH//PD\). 而 \(PD\subset\) 平面 \(PAD\)，且 \(GH\) 不在平面 \(PAD\) 内，所以 \(GH//\) 平面 \(PAD\).

\item 取 \(PC\) 的中点 \(N\)，连接 \(DN\). 因为 \(\triangle PCD\) 为等边三角形，所以 \(DN\perp PC\). 又平面 \(PAC\perp\) 平面 \(PCD\)，且两平面的交线为 \(PC\)，所以 \(DN\perp\) 平面 \(PAC\). 从而 \(DN\perp PA\). 又 \(PA\perp CD\)，且 \(DN\subset\) 平面 \(PCD\)，所以 \(PA\) 同时垂直于平面 \(PCD\) 内的两条相交直线 \(CD\) 与 \(DN\)，故 \(PA\perp\) 平面 \(PCD\).

\item 连接 \(AN\). 由（2）知 \(DN\perp\) 平面 \(PAC\)，所以 \(\angle DAN\) 是直线 \(AD\) 与平面 \(PAC\) 所成的角. 因为 \(\triangle PCD\) 为等边三角形，\(CD=2\)，且 \(N\) 为 \(PC\) 的中点，所以 \(DN=\sqrt3\). 又 \(AD=3\)，故
\(\sin\angle DAN=\dfrac{DN}{AD}=\dfrac{\sqrt3}{3}\).
\end{enumerate}
\end{solution}

\begin{problem}
设 \(\{a_n\}\) 是等差数列，\(\{b_n\}\) 是等比数列，公比大于 \(0\). 已知 \(a_1=b_1=3\)，\(b_2=a_3\)，\(b_3=4a_2+3\).

\begin{enumerate}
\item 求 \(\{a_n\}\) 和 \(\{b_n\}\) 的通项公式；
\item 设数列 \(\{c_n\}\) 满足：当 \(n\) 为奇数时，\(c_n=1\)；当 \(n\) 为偶数时，\(c_n=b_{n/2}\). 求 \(a_1c_1+a_2c_2+\cdots+a_{2n}c_{2n}\)（\(n\in\N^*\)）.
\end{enumerate}
\end{problem}

\begin{solution}
\begin{enumerate}
\item 设等差数列 \(\{a_n\}\) 的公差为 \(d\)，等比数列 \(\{b_n\}\) 的公比为 \(q\). 由 \(a_1=b_1=3\)，\(b_2=a_3\)，\(b_3=4a_2+3\)，得 \(3q=3+2d\)，\(3q^2=15+4d\). 解得 \(d=3\)，\(q=3\). 所以 \(a_n=3+3(n-1)=3n\)，\(b_n=3\cdot 3^{n-1}=3^n\).

\item
\(a_1c_1+a_2c_2+\cdots+a_{2n}c_{2n}=\bigl(a_1+a_3+\cdots+a_{2n-1}\bigr)+\bigl(a_2b_1+a_4b_2+\cdots+a_{2n}b_n\bigr)\).
其中
\(a_1+a_3+\cdots+a_{2n-1}=3+9+\cdots+(6n-3)=3n^2\)，
又
\(a_2b_1+a_4b_2+\cdots+a_{2n}b_n=6(3+2\cdot 3^2+\cdots+n\cdot 3^n)\).
设
\(T_n=1\cdot 3^1+2\cdot 3^2+\cdots+n\cdot 3^n\)，
则
\(3T_n=1\cdot 3^2+2\cdot 3^3+\cdots+n\cdot 3^{n+1}\).
两式相减得
\(2T_n=-(3+3^2+\cdots+3^n)+n\cdot 3^{n+1}=-\dfrac{3(1-3^n)}{1-3}+n\cdot 3^{n+1}\).
整理得
\(T_n=\dfrac{(2n-1)3^{n+1}+3}{2}\).
故
\(a_1c_1+a_2c_2+\cdots+a_{2n}c_{2n}=3n^2+6T_n=\dfrac{(2n-1)3^{n+2}+6n^2+9}{2}\).
\end{enumerate}
\end{solution}

\begin{problem}
设椭圆 \(\dfrac{x^2}{a^2}+\dfrac{y^2}{b^2}=1\)（\(a>b>0\)） 的左焦点为 \(F\)，左顶点为 \(A\)，上顶点为 \(B\). 已知 \(\sqrt3|OA|=2|OB|\)（\(O\) 为原点）.

\begin{enumerate}
\item 求椭圆的离心率；
\item 设经过点 \(F\) 且斜率为 \(\dfrac34\) 的直线 \(l\) 与椭圆在 \(x\) 轴上方的交点为 \(P\)，圆 \(C\) 同时与 \(x\) 轴和直线 \(l\) 相切，圆心 \(C\) 在直线 \(x=4\) 上，且 \(OC//AP\)，求椭圆的方程.
\end{enumerate}
\end{problem}

\begin{solution}
\begin{enumerate}
\item 设椭圆的半焦距为 \(c\). 由 \(\sqrt3|OA|=2|OB|\) 可得 \(\sqrt3\,a=2b\). 又 \(c^2=a^2-b^2\)，消去 \(b\) 得 \(a=2c\). 所以椭圆的离心率 \(e=\dfrac ca=\dfrac12\).

\item 由（1）知 \(a=2c\)，\(b=\sqrt3c\)，故椭圆方程为 \(\dfrac{x^2}{4c^2}+\dfrac{y^2}{3c^2}=1\). 又 \(F(-c,0)\)，故直线 \(l\) 的方程为 \(y=\dfrac34(x+c)\). 联立可得点 \(P\) 的横坐标满足 \(7x^2+6cx-13c^2=0\)，于是 \(x_1=c\)，\(x_2=-\dfrac{13c}{7}\). 代入直线 \(l\) 的方程，得 \(y_1=\dfrac32c\)，\(y_2=-\dfrac9{14}c\). 由于 \(P\) 在 \(x\) 轴上方，故取 \(P\left(c,\dfrac32c\right)\). 设圆心 \(C(4,t)\). 由 \(OC//AP\) 得 \(\dfrac t4=\dfrac{\frac32c}{c+2c}=\dfrac12\)，故 \(t=2\). 又圆 \(C\) 与 \(x\) 轴相切，所以半径为 \(2\). 再由圆 \(C\) 与 \(l\) 相切，可得 \(\dfrac{\left|\frac34(4+c)-2\right|}{\sqrt{1+\left(\frac34\right)^2}}=2\)，解得 \(c=2\). 于是椭圆方程为 \(\dfrac{x^2}{16}+\dfrac{y^2}{12}=1\).
\end{enumerate}
\end{solution}

\begin{problem}
设函数 \(f(x)=\ln x-a(x-1)\e^x\)，其中 \(a\in\R\).

\begin{enumerate}
\item 若 \(a\le0\)，讨论 \(f(x)\) 的单调性；
\item 若 \(0<a<\dfrac1\e\)，
\begin{enumerate}
\item 证明 \(f(x)\) 恰有两个零点；
\item 设 \(x_0\) 为 \(f(x)\) 的极值点，\(x_1\) 为 \(f(x)\) 的零点，且 \(x_1>x_0\)，证明 \(3x_0-x_1>2\).
\end{enumerate}
\end{enumerate}
\end{problem}

\begin{solution}
\begin{enumerate}
\item \(f(x)\) 的定义域为 \((0,+\infty)\)，且 \(f'(x)=\dfrac1x-a x\e^x\). 当 \(a\le0\) 时，\(f'(x)>0\)，故 \(f(x)\) 在 \((0,+\infty)\) 内单调递增.

\item
\begin{enumerate}
\item 令 \(g(x)=1-a x^2\e^x\). 则 \(f'(x)=\dfrac{g(x)}{x}\). 由于 \(0<a<\dfrac1\e\)，故 \(g(x)\) 在 \((0,+\infty)\) 内单调递减，且 \(g(1)=1-a\e>0\)，\(g\!\left(\ln\dfrac1a\right)=1-a\left(\ln\dfrac1a\right)^2\cdot\dfrac1a<0\). 所以 \(g(x)=0\) 在 \((0,+\infty)\) 内有唯一解，不妨记为 \(x_0\)，且 \(1<x_0<\ln\dfrac1a\). 于是 \(f(x)\) 在 \((0,x_0)\) 内单调递增，在 \((x_0,+\infty)\) 内单调递减，因此 \(x_0\) 是唯一极值点.
又令 \(h(x)=\ln x-x+1\)，则当 \(x>1\) 时，\(h'(x)=\dfrac1x-1<0\)，故 \(h(x)\) 在 \((1,+\infty)\) 内单调递减，并且 \(h(1)=0\). 从而当 \(x>1\) 时，\(\ln x<x-1\). 又 \(f\!\left(\ln\dfrac1a\right)=\ln\ln\dfrac1a-a\left(\ln\dfrac1a-1\right)\cdot\dfrac1a=\ln\ln\dfrac1a-\ln\dfrac1a+1=h\!\left(\ln\dfrac1a\right)<0\). 因为 \(f(x_0)>f(1)=0\)，所以 \(f(x)\) 在 \((x_0,+\infty)\) 内有唯一零点. 又 \(1<x_0\)，且 \(f(1)=0\)，所以 \(f(x)\) 在 \((0,x_0)\) 内有唯一零点 \(1\). 因此 \(f(x)\) 恰有两个零点.

\item 设 \(x_0\) 为 \(f(x)\) 的极值点，\(x_1\) 为 \(f(x)\) 的零点，且 \(x_1>x_0\). 由题意，\(f'(x_0)=0\)，\(f(x_1)=0\). 即 \(\dfrac1{x_0}=a x_0\e^{x_0}\)，\(\ln x_1=a(x_1-1)\e^{x_1}\)，从而 \(\ln x_1=\dfrac{x_1-1}{x_0^2}\e^{x_1-x_0}\). 由 \(x_1>x_0>1\) 可得 \(\ln x_1<x_1-1\)，从而 \(\e^{x_1-x_0}<x_0^2\). 两边取对数得 \(x_1-x_0<2\ln x_0<2(x_0-1)\)，于是 \(3x_0-x_1>2\).
\end{enumerate}
\end{enumerate}
\end{solution}
